\chapter{2012 Nobel Prize in Economics}
\textbf{Laureates: }Alvin Roth (USA), Lloyd Shapley (USA)

\section*{Reason for Award}
For the theory of stable allocations and the practice of market design

\section{What They Did}
Lloyd Shapley and Alvin Roth solved the stable matching problem, developing
algorithms that produce stable assignments in two-sided markets. Shapley proved
the existence of stable matchings and developed the deferred acceptance algorithm,
which guarantees stability regardless of preference structures.

His work showed that the algorithm produces a matching that is optimal for one
side of the market while being strategy-proof for the other side. Roth and
Shapley showed how this theory applies to practical market design problems.

They analyzed matching markets for medical residents, school choice programs,
and kidney exchanges. Roth identified when markets fail and require design
intervention, developing the concept of manipulability and proving that some
matching mechanisms cannot be made strategy-proof.

Shapley's earlier work on cooperative game theory and the Shapley value provided
the theoretical foundation for understanding fair allocation in collaborative
settings. The Shapley value assigns each player a payoff proportional to their
marginal contribution across all possible coalitions, providing a principled
method for dividing surplus among participants.

Roth recognized that the matching market for medical residents was failing.
The original match, designed by Shapley and others, had been manipulated by
hospitals, making it unstable and inefficient. Roth redesigned the mechanism,
creating the National Resident Matching Program (NRMP) algorithm that has
successfully matched over 30,000 medical graduates to residency positions
since 1998.

The theory also applies to school choice systems. Roth and co-authors designed
the Boston Public Schools match, which initially had problems with strategic
manipulation. After criticism from economists, the system was redesigned to use
the deferred acceptance algorithm, dramatically improving fairness and efficiency.

\section{Impact on Economic Thinking}
Shapley and Roth transformed understanding of market design and matching theory.
Their work showed that markets are not natural phenomena but human constructions
that can be designed for better outcomes. The stable matching framework provided
tools for analyzing allocation problems where prices are prohibited or impractical.

This influenced public policy regarding school choice, organ transplantation,
and housing allocation. Roth's practical interventions improved matching
mechanisms worldwide, demonstrating that economic theory could directly enhance
social welfare. The redesign of the Boston school choice system showed that
theoretical insights could be rapidly applied to improve real-world outcomes.

The work challenged the neoclassical assumption that markets always clear through
price adjustments, showing that many important allocations occur without prices.
Labor markets, housing markets, and educational placements often cannot use prices
due to legal, ethical, or practical constraints. Stable matching provides the
theoretical foundation for these non-price allocation mechanisms.

Their theory also influenced computer science, game theory, and operations
research, creating interdisciplinary connections that enriched all fields. The
algorithmic foundations they established provided computational tools for
solving large-scale matching problems in practice, with applications ranging
from kidney exchange to spectrum allocation.

The concept of stability has become central to mechanism design theory, providing
a benchmark for evaluating the quality of allocation mechanisms. Researchers
now routinely analyze whether proposed mechanisms produce stable outcomes and
how they perform under strategic behavior.

\section{Modern Perspectives Today}
Market design based on Shapley and Roth's work has expanded to diverse
applications. School choice systems in New York, Boston, and other cities use
stable matching algorithms to assign students to schools, reducing stratification
and improving access to quality education.

Kidney exchange programs save lives by matching incompatible donor-patient pairs
across hospitals. The National Kidney Registry in the United States uses algorithms
based on stable matching to facilitate over 5,000 transplants annually. Roth has
advocated for expanded use of these programs to address the chronic shortage of
kidneys for transplant.

Spectrum auctions, blood donation systems, and even dating algorithms draw on
matching theory. The Federal Communications Commission has used auction designs
informed by matching theory to allocate radio frequencies efficiently.

Roth has advocated for market design solutions to organ shortage crises,
proposing paired exchange and priority markets. The theory has been extended to
dynamic matching, heterogeneous agents, and network-based allocations.

Digital platforms like ride-sharing and food delivery use matching algorithms
that embody stable matching principles. Uber, Lyft, and DoorDash employ
real-time matching algorithms that balance supply and demand across geographic
spaces. Climate policy mechanisms, such as carbon credit trading, also apply
market design insights to allocate emission permits efficiently.

As matching theory continues to evolve, its applications expand to emerging
challenges like organ allocation during pandemics, school choice in the age of
virtual learning, and resource allocation in digital marketplaces.
